\subsection{Historical approaches}
\begin{quotation}\textit{
The art of cryptography is all about trying to make it difficult for
your adversary to decipher the secret message you have encoded. Going back
to classical times we have the Caesar Cipher, where each letter in the
plaintext is replaces by the letter some fixed number of places beyond it
in the alphabet. This is of course amazingly easy to break by just trying
a range of offsets. Puzzle-makers enjoy finding words or even phrases that
are readable in both origibnal and shifted form. For instance ``ferns'' and
``banjo'' are alphabetically shifted relative to each other by 4. So consider
finding the longest or most interesting such example as a challenge!
Over the centuries many (much) stronger coding schemes were invented. Perhaps
the culmination of what can be characterised as ``historical'' was in
schemes used during the 1939-45 war where serious energy went into
proto-computer schemes to break the encryption.
}\end{quotation}
\subsubsection{The one-time pad}
\subsubsection{Fish}
\subsubsection{Enigma}

\subsection{Modern Cryptography}
\begin{quotation}\textit{
Once computers started to be used for code-breaking there was an
obvious need to use them for encoding. Two particular schemes stand out
as both easy to describe and as widely used.
}\end{quotation}
\subsubsection{The alleged RC4 stream cipher}
\subsubsection{Public key and RSA}

\subsection{Looking forward to a post-quantum world}
\begin{quotation}\textit{
At the time of writing there is a serious concern that future large scale
computers exploiting quantum rather than classical physics might be able
to break pretty well all existing encryption schemes. This has led to
work to develop methods that should remain robustly secure even in the
antifipated new world, and a gradual transition away from using traditional
methods to application of these new quantum-resistent techniques. The
technicalities behind all of this are fairly severe, so the discussion
here is going to be rather abbreviated!
}\end{quotation}

\subsection{Keeping secrets os hard}
Beyond the details of how to encrypt data in a way that is secure not just
today but in the future, the protocols for the exchange of secrets can
be hard to get right. And changes in computer technology (such as first the
availability of massive parallel computer farms, and perhaps soon
quantum computers) means that one can not afford to stand still. It is
a brutally difficult world to work in. 
